% 1- Explicar detenidamente la principal motivacion del articulo. Dejar claro cual es el posible problema
% 2- Presentar las preguntas de investigacion pertinentes
% 3- Presentar algunos resultados preliminares
% 4- Presentar el trabajo llevado hasta ahora de investigacion y desarrollo
% 5- Dejar algunas preguntas abiertas y problemas encontrados hasta ahora
%
% $Id: slides.tex 4228 2006-06-21 21:55:12Z jjamor $
%
%
% Compilar a .pdf con LaTeX (pdflatex)
% Es necesario instalar Beamer (paquete latex-beamer en Debian)
%

%
% Grficos:
% Los grficos pueden suministrarse en PNG, JPG, TIF, PDF, MPS
% Los EPS deben convertirse a PDF (usar epstopdf)
%


\documentclass{beamer}
\usetheme{Warsaw}
\usebackgroundtemplate{\includegraphics[width=\paperwidth]{format/libresoft-bg.png}}
\usepackage[english]{babel}
\usepackage[latin1]{inputenc}
\usepackage{graphics}
\usepackage{amssymb} % Simbolos matematicos
\usepackage{url}

%\definecolor{libresoftgreen}{RGB}{162,190,43}
%\definecolor{libresoftblue}{RGB}{0,98,143}

%\setbeamercolor{titlelike}{bg=libresoftgreen}

%% Metadatos del PDF.
\hypersetup{
  pdftitle={Libresoft},
  pdfauthor={Daniel Izquierdo Cort\'azar},
  pdfcreator={GSyC/Libresoft},
  pdfproducer=PDFLaTeX,
  pdfsubject={QualOSS: Quality of Open Source Software},
}
%%



\AtBeginSection[]
{
  \begin{frame}<presentation>
    \frametitle{Index}
    \tableofcontents[current]
  \end{frame}
}




\begin{document}



\title{QualOSS: Quality of Open Source Software
}
\subtitle{ Master on Free Software
}
\institute{dizquierdo@gsyc.es\\
GSyC/Libresoft, Universidad Rey Juan Carlos}
\author{Felipe Ortega, Pedro Coca, Daniel Izquierdo Cort\'azar}
\date{Fuenlabrada, Madrid\\ 2nd, December 2010}

\frame{
\maketitle
\begin{center}
\includegraphics[width=4cm]{format/gsyc-urjc.pdf} 
\end{center}
}


% Si el titulo o el autor se quieren acortar para los pies de pgina
% se pueden redefinir aqu:
%\title{Titulo corto}
%\author{Autores abreviado}


%% LICENCIA DE REDISTRIBUCION DE LAS TRANSPAS
\frame{
~
\vspace{4cm}

\begin{flushright}
{\tiny
(cc) 2010 Daniel Izquierdo Cort\'azar. \\
Some rights reserved. This document is distributed under the Creative \\
            Commons Attribution-ShareAlike 2.5 licence, available in \\
            http://creativecommons.org/licenses/by-sa/2.5/

%  Este documento (o uno muy similar) está disponible en \\
%  \url{http://gsyc.escet.urjc.es/~jjamor/}
}
\end{flushright}
}
%%

\frame{
\frametitle{Table of contents}
\tableofcontents
}


\section{Introduction}


\begin{frame}
\frametitle{Motivation}
\begin{center}
\begin{itemize}
\item Plenty of companies use FLOSS products.
\item And when acquiring software, they are not only interested to know about the product,
\item but also about who produced that and how reliable they are.

\end{itemize}
\end{center}
\end{frame}



\begin{frame}
\frametitle{QualOSS}
\begin{center}
\begin{itemize}
\item Not too many quality models are focused on FLOSS so far
\item and if so, they are still too light weight : OpenBRR / QSoS

\end{itemize}
\end{center}
\end{frame}



\begin{frame}
\frametitle{QualOSS}
\begin{center}
\begin{itemize}
\item QualOSS: Quality of Open Source Software
\item http://qualoss.eu
\item This project aimed to fill that gap in the state of the art

\end{itemize}
\end{center}
\end{frame}




\begin{frame}
\frametitle{Automation}
\begin{center}
\begin{itemize}
\item One of its goals consisted of semi-full automation of the analysis
\item The community metrics were retrieved using FLOSSMetrics (http://flossmetrics.org) data.
\item And several scripts are provided for this purpose.
\end{itemize}
\end{center}
\end{frame}


\section{Methodology}
\subsection{Initial Steps}

\begin{frame}
\frametitle{Initial set of potential quality attributes}
\begin{center}
\begin{itemize}
\item Interviews with several companies.
\item Some of them were interested in communities.
\item In fact, some of them were just interested in the community and not the source code.
\item Further information can be found at deliverable 1.2
\end{itemize}
\end{center}
\end{frame}



\begin{frame}
\frametitle{Some of the questions}
\begin{center}
\begin{itemize}
\item\textit{ For which types of purposes/systems do you (would you) use F/OSS products?}
\item \textit{What are criteria that F/OSS components need to fulfill to be appropriate?}
\item\textit{ To your knowledge, do you achieve business goals through F/OSS usage?}
\end{itemize}
\end{center}
\end{frame}


\begin{frame}
\frametitle{Initial set of potential quality attributes}
\begin{center}
\includegraphics[width=0.9\textwidth]{figs/initialqualityattributes.png}
\end{center}
\end{frame}

\begin{frame}
\frametitle{Initial set of potential quality attributes}
\begin{center}
\begin{itemize}
\item One of the main attributes raised during the interviews was
the idea of measuring community.
\begin{itemize}
\item Is able a community to survive for years?
\item Imagine that the main developer of a community disappears, how does this affect the evolvability of the product?
\end{itemize}
\end{itemize}
\end{center}
\end{frame}

\subsection{Basic Quality Model: Methodology}

\begin{frame}
\frametitle{Main quality attributes}
\begin{center}
\begin{itemize}
\item The two key quality attributes were Robustness and Evolvability
\item For each of them it was followed a top-down approach
\end{itemize}
\end{center}
\end{frame}


\begin{frame}
\frametitle{Goal - Question - Metric}
\begin{center}
\begin{itemize}
\item It follows the GQM approach.
\item Each goal (evolvability and robustness) is divided into several questions.
\item And several metrics are defined for each question.

\end{itemize}
\end{center}
\end{frame}


\begin{frame}
\frametitle{QualOSS Quality Model}
\begin{center}
\begin{itemize}
\end{itemize}
\includegraphics[width=0.9\textwidth]{figs/rob_evol_4StdQualossMethod.png}
\end{center}
\end{frame}


\begin{frame}
\frametitle{QualOSS Quality Model}
\begin{center}
\begin{itemize}
\item Let's focus on the community side
\end{itemize}

\end{center}
\end{frame}


\subsection{Community Side}

\begin{frame}
\frametitle{Definitions of the main quality attributes}
\begin{center}
\begin{itemize}
\item Size and Regeneration Adequacy: The degree to which the size evolution and regeneration of a FLOSS community happens at an adequate rate to maintain a sustainable community size.

\item Interactivity and Workload Adequacy:   The degree to which the community interacts adequately and partition the workload among FLOSS community members adequately to maintain a community cohesion and motivated.
\end{itemize}
\end{center}
\end{frame}


\begin{frame}
\frametitle{Questions: Size and Regeneration Adequacy}
\begin{center}
\begin{itemize}
\item Has the evolution of new community members reporting bugs re-
mained stable or grown over the history of the FLOSS endeavor? (at
least shown a stable or positive trend overall)
\item  Has the evolution of new code contributing members remained stable
or grown over the history of the FLOSS endeavor? (at least shown a
stable or positive trend overall)
\item  Has the evolution of new members contributing data other than code
or bug report remained stable or grown over the history of the FLOSS
endeavor? (at least shown a stable or positive trend overall)
\end{itemize}
\end{center}
\end{frame}

\begin{frame}
\frametitle{Questions: Size and Regeneration Adequacy}
\begin{center}
\begin{itemize}
\item  Has the evolution of new core contributing members remained stable
or grown over the history of the FLOSS endeavor? (at least shown a
stable or positive trend overall) (a core member is one with commit
right who perform commits frequently for instance, more than once
every three month period)
\item  What is the evolution of core members who stopped contributing for
a significant period?
\item  Has the evolution of core members who stopped contributing for a
significant period been compensated by the joining of new core members around the same time frame?
\item What is the average longevity of committers to the FLOSS endeavor

\end{itemize}
\end{center}
\end{frame}


\begin{frame}
\frametitle{Questions: Community Interactivity and Workload Adequacy}
\begin{center}
\begin{itemize}
\item  Is the number of events adequate (to show a lively community)?
\item Is the number of code commits adequate (to show a lively committer
community)?
\item Has the size of community supporting older versions of a FLOSS
component remained sufficient? (compared to number of bug reports,
the number of report fix vs those open and compared to code size)
\item Are they sub-groups in the community? If so, are they disconnected
or are active community members serving as bridges between these
sub groups (sub groups can be at the level of roles such bug reporter,
committer but subgroups can also be studied at the code level)

\end{itemize}
\end{center}
\end{frame}





\begin{frame}
\frametitle{Questions: Community Composition Adequacy}
\begin{center}
\begin{itemize}
\item  What companies have contributed bug reports?
\item What companies have contributed code (code patches or code com-
mits)?
\item What companies have employees in leadership position in the FLOSS
community (project leader, release manager, treasurer, ...)?
\item What companies provide services on the FLOSS component?
\item What are the roles who are filled by active community members?
\item Are these roles adequate for the FLOSS endeavor in question?
\item Is there a sufficient number of code expert on the various portion of
the FLOSS component?


\end{itemize}
\end{center}
\end{frame}


\subsection{Metrics}

\begin{frame}
\frametitle{Metrics}
\begin{center}
\begin{itemize}
\item They should be as objective as possible
\item One of more metrics are matched with a given question.
\item \textbf{S}ize and \textbf{R}egeneration \textbf{A}dequacy (SRA)
\item  Community \textbf{I}nteractivity and \textbf{W}orkload \textbf{A}dequacy (IWA)
\end{itemize}
\end{center}
\end{frame}



\begin{frame}
\frametitle{Metrics: Averages}
\begin{center}
\begin{itemize}
\item Metrics whose results is directly calculated as an average among some data:
\begin{itemize}
\item \textbf{sra7} average number of months where each committer committed,
\item \textbf{iwa4} proportion of files maintained by a single committer,
\item \textbf{iwa5} ratio of the number of lines of source code by the number of committers
active in the last year,
\item \textbf{iwa7} proportion of code files committed to in the last year.
\end{itemize}

\end{itemize}
\end{center}
\end{frame}



\begin{frame}
\frametitle{Metrics: Describing evolution of the community}
\begin{center}
\begin{itemize}
\item Metrics whose results aim to represent the evolution of a given concept:
\begin{itemize}
\item \textbf{sra2} number of new code committers in the interval,
\item \textbf{sra3} number of new non code committers in the interval,
\item \textbf{sra9} number of active code committers in the interval,
\item \textbf{iwa1} ratio of the number of commits by the number of committers in the
interval,
\item \textbf{iwa2} ratio of the number of code commits by the number of code committers
in the interval.
\end{itemize}
\end{itemize}
\end{center}
\end{frame}



\begin{frame}
\frametitle{Metrics: Describing the evolution of the community}
\begin{center}
\begin{itemize}
\item Next metrics use the notion of \textit{core committer}
\begin{itemize}
\item \textbf{sra4} number of new core committers in the year
\item \textbf{sra5} number of core committers that quit in the year
\item \textbf{sra6} difference between sra4 and sra5.

\end{itemize}
\end{itemize}
\end{center}
\end{frame}



\subsection{Obtaining Metrics}

\begin{frame}
\frametitle{Methodology}
\begin{center}
\begin{itemize}
\item To measure risky situations, thresholds were defined: green, yellow, red and black.
\item Those thresholds were derived from a statistical analysis of a set of 1400 projects from FLOSSMetrics dataset.
\item Those thresholds show how a given project is located among others. This could be useful from
the point of view of benchmarking. 
\item Having a black color does not mean that this project is ``bad'', but it means that the results were in the set of projects that are most far from the average.

\end{itemize}
\end{center}
\end{frame}





\section{Results}


\begin{frame}
\frametitle{Some Results}
\begin{center}
\begin{itemize}
\item sra7: Average longevity measured in months.
\item iwa4: Territoriality.
\item sra2: New code core contributors.\\
\end{itemize}

\vspace{0.5cm}

\begin{tabular}{l*{4}{c}}
%\toprule
Pred.\ diff.  & evolution & evince & httpd1.3\\
\hline\noalign{\smallskip}

sra7 &  8.6044 Y   & 5.0272 B   & 15.7206 G\\
\hline\noalign{\smallskip}
iwa4 &  0.3666 G   & 0.4216 G   & 0.0923 G\\
%iwa5 &  374772.4 B & 48282.8 R  & -\\
%iwa7 &  0.145 G    & 0.1904 Y   & 6.0e-4 B\\
\hline\noalign{\smallskip}
%$\Delta_{90}$ sra2$_30}   & 0 Y & 0 Y & 0 Y & 0 Y\\
$\Delta_{90}$ sra2$_{90}$   & 0 Y & 0 Y & 0 Y\\
%$\Delta_{90}$ sra2$_180}  & 0 Y & 0 Y & 0 Y & 0 Y\\
%$\Delta_{180}$ sra2$_30}  & 0 Y & 0 Y & 0 Y & 0 Y\\
%$\Delta_{180}$ sra2$_90}  & 0 Y & 0 Y & 0 Y & 0 Y\\
%$\Delta_{180}$ sra2$_{180}$ & 0 Y & 0 Y & 0 Y & 0 Y\\

%$\Delta_{90}$ sra3$_30}   & 0 Y & 0 Y & 0 Y & 0 Y\\
   %$\Delta_{90}$ sra3$_{90}$   & 0 Y & 0 Y & 0 Y\\
%$\Delta_{90}$ sra3$_180$  & 0 Y & 0 Y & 0 Y & 0 Y\\
%$\Delta_{180}$ sra3$_30}  & 0 Y & 0 Y & 0 Y & 0 Y\\
%$\Delta_{180}$ sra3$_90}  & 0 Y & 0 Y & 0 Y & 0 Y\\
%$\Delta_{180}$ sra3$_{180}$ & 0 Y & 1 Y & 0 Y & 0 Y\\

%$\Delta_{90}$ sra9$_30}   & 0 Y & 0 Y & 0 Y & 0 Y\\
   %$\Delta_{90}$ sra9$_{90}$   & 0 Y & 0 Y & 0 Y\\
%$\Delta_{90}$ sra9$_180}  & 0 Y & 0 Y & 0 Y & 0 Y\\
%$\Delta_{180}$ sra9$_30}  & 0 Y & 0 Y & 0 Y & 0 Y\\
%$\Delta_{180}$ sra9$_90}  & 0 Y & 0 Y & 0 Y & -1 R\\
%$\Delta_{180}$ sra9$_{180}$ & 0 Y & 0 Y & 0 Y & -1 R\\

%$\Delta_{90}$ iwa1$_30}  & -0.11461449 R & 0.01739115 Y & -0.09101181 R & -0.26294442 R\\
    %$\Delta_{90}$ iwa1$_{90}$  & -0.23327969 R & -0.0312513 R & -0.62113349 R\\
%$\Delta_{90}$ iwa1$_180} & -0.308871655 R & -0.12249361 R & -0.36510017 R & -1.09633002R\\
%$\Delta_{180}$ iwa1$_30} & -0.22922898 Y & 0.0347823 Y & -0.18202362 Y & -0.52588884 Y\\
%$\Delta_{180}$ iwa1$_90} & -0.46655938 Y & -0.0625026 Y & -0.40314032 Y & -1.24226698 R\\
%$\Delta_{180}$ iwa1$_{180}$& -0.61774331 Y & -0.24498722 Y & -0.73020034 Y & -2.19266004 R\\

%$\Delta_{90}$ iwa2$_30}  & -0.53154186 R & 0.37025139 Y & -0.79931637 R & -0.26808226 R\\
     %$\Delta_{90}$ iwa2$_{90}$  & -1.03448107 R & 0.31998312 Y & -0.76343814 R\\
%$\Delta_{90}$ iwa2$_180} & -1.236808695 R & -0.078177255 R & -2.75218718 R & -1.3574936 R\\
%$\Delta_{180}$ iwa2$_30}  & -1.06308372 Y & 0.74050278 Y & -1.59863274 Y & -0.5361645 Y\\
%$\Delta_{180}$ iwa2$_90}  & -2.06896214 R & 0.63996624 Y & -3.21783906 R & -1.52687628 Y\\
%$\Delta_{180}$ iwa2$_{180}$ & -2.47361739 R & -0.15635451 Y & -5.50437436 R & -2.71498719 R\\

%\hline\noalign{\smallskip}
%$\Delta$ sra4  & 1 G & 2 G &  0 Y\\
%$\Delta$ sra5 & 1 B & 1 B &  0 R\\
%$\Delta$ sra6 & -1 R & 0 Y &  0 Y\\
%\hline\noalign{\smallskip}
\end{tabular}

\end{center}
\end{frame}


\begin{frame}
\frametitle{Results}
\begin{center}
\begin{itemize}
\item At CETIC website: \url{http://ingrid.cetic.be:33323/qualoss_assessment/}
\item At FLOSSMetrics website: \url{http://melquiades.flossmetrics.org/projects/octave/quality}
\end{itemize}
\end{center}
\end{frame}






\begin{frame}
\frametitle{Conclusions}
\begin{center}
\begin{itemize}
\item Statistical approach to create a benchmarking to give a mark for a given project and for a given metric.
\item This could be useful to check the tendency of the community since part of the metrics measure the evolution of some characteristics.
\item As said, indicators measure relative and not absolute risk. 
\item 22 metrics proposed for the community side, where 16 of them are fully automated using data 
from FLOSSMetrics.
\end{itemize}
\end{center}
\end{frame}

%**********************


\begin{frame}
\frametitle{Questions?}
\begin{center}
%Thanks for your attendance!\\
Questions?\\
%--\\
%\begin{small}
%Daniel Izquierdo Cortzar\\
%dizquierdo@gsyc.es\\
%GSyC/LibreSoft - Universidad Rey Juan Carlos\\
%\end{small}
\end{center}
\end{frame}




\end{document}
